\documentclass[11pt,a4paper]{article}

% --- Mathematics (must precede unicode-math) ---
\usepackage{amsmath}
\usepackage{mathtools}

% --- Fonts (lualatex) ---
\usepackage{fontspec}
\usepackage{unicode-math}
\setmathfont{KpMath-Regular.otf}
\setmathfont{KpMath-Bold.otf}[version=bold]

\usepackage{microtype}
\usepackage[dvipsnames]{xcolor}
\usepackage[margin=25mm, top=30mm, bottom=30mm]{geometry}

\usepackage{fancyhdr}
\pagestyle{fancy}
\fancyhf{}
\fancyhead[L]{\small\itshape The Continuum Limit of a Space-Filling Layer of Cubes}
\fancyhead[R]{\small\itshape Nestor (2026)}
\fancyfoot[C]{\thepage}
\renewcommand{\headrulewidth}{0.4pt}
\renewcommand{\footrulewidth}{0pt}

\setcounter{secnumdepth}{3}
\usepackage{booktabs}
\usepackage{array}
\usepackage[font=small, labelfont=bf, labelsep=period]{caption}
\usepackage{setspace}
\setstretch{1.15}
\setlength{\parskip}{0.4em}
\usepackage[colorlinks=true, linkcolor=blue!60!black,
            citecolor=blue!60!black, urlcolor=blue!60!black]{hyperref}

% --- Macros ---
\newcommand{\bI}{\mathbf{I}}
\newcommand{\bA}{\mathbf{A}}
\newcommand{\bS}{\mathbf{S}}
\newcommand{\bT}{\mathbf{T}}
\newcommand{\bK}{\mathbf{K}}
\newcommand{\bP}{\mathbf{P}}
\newcommand{\bD}{\boldsymbol{\Delta}}
\newcommand{\bG}{\mathbf{G}}
\newcommand{\sinc}{\operatorname{sinc}}
\newcommand{\nb}[1]{\texttt{#1}}

\title{\bfseries The Continuum Limit of a Space-Filling Layer of Cubes\\[4pt]
\large A Tiling Identity, an Exact Lattice Kernel,\\
and the Closed-Form Error of the Discrete Layer}
\author{}
\date{28 September 2026}

\begin{document}
\maketitle

\begin{abstract}
\noindent
A heterogeneous elastic medium discretised into cubic voxels, each carrying a single-site
$T$-matrix and coupled through a lattice Green's tensor, must reproduce the continuous medium as
the voxels shrink. For a laterally uniform layer at normal incidence we establish when, how fast,
and with what error it does. Five results. (i) Cubes tile the plane, so the lateral form factor of
a cube vanishes on the reciprocal lattice; consequently the source-cell-averaged coupling between
planes is \emph{exactly} the continuum's plane-wave term, and within a plane the sum over all
cells, the voxel's own included, is \emph{exactly} the field of a uniform plate. This tiling
identity gives the lattice kernel at normal incidence in closed form, with no lattice sum.
(ii) The collocation closure of the single site carries the same self term the kernel removes, so
the self term cancels and the discrete system becomes a collocation of the continuum equation with
exact cell integrals. (iii) That collocation converges at exactly second order in the voxel size.
(iv) Its leading error is closed-form for every discretisation --- a midpoint-rule error of the
internal field across each voxel, $+(kd)^2/8$ in reflection and $-(kd)^2/24$ in transmission ---
and the opposite signs show that no scalar correction of a uniform-field single site can remove it.
(v) A voxel that carries the first moment of its field, coupled through exact cell-to-cell integrals,
removes it and converges at fourth order: one plane of such voxels is $6500$ times more accurate
than one plane of uniform ones, and the scheme stays fourth order at oblique incidence, in the
converted S wave as well as the reflected P. Every result is derived symbolically in
Mathematica and checked against an independent numerical implementation.
\end{abstract}

\section{Setting}\label{sec:setting}

\paragraph{The discrete medium.} Space is tiled by cubes of side $d$ (half-side $h=d/2$). Cube $i$
carries a contrast $\Delta\lambda,\Delta\mu,\Delta\rho$ from an isotropic background with Lamé
parameters $\lambda,\mu$, density $\rho$, P and S speeds $\alpha,\beta$ and P-wave modulus
$M_P=\lambda+2\mu$. Time dependence is $e^{-i\omega t}$; $k_P=\omega/\alpha$, $k_S=\omega/\beta$.
Coordinates are ordered $(z,x,y)$ with $z$ down. The field at a voxel is the nine-component state
\begin{equation}
  \psi=(u_z,u_x,u_y,\ \varepsilon_{zz},\varepsilon_{xx},\varepsilon_{yy},
  2\varepsilon_{xy},2\varepsilon_{zy},2\varepsilon_{zx}),
\end{equation}
displacement followed by engineering strain. A voxel responds with a force and a moment (stress
dipole); the $9\times9$ single-site $T$-matrix maps the exciting state to them, and the $9\times9$
lattice Green's tensor $\bK(i,j)$ maps a source voxel's force and moment to the state at voxel $i$.
The Foldy--Lax system is
\begin{equation}\label{eq:foldylax}
  \psi_i=\psi^0_i+\sum_{j\neq i}\bK(i,j)\,\bT_j\,\psi_j ,
\end{equation}
with $\psi^0$ the incident state. Strain rows of $\bK$ are receiver derivatives of the point
Green's tensor $\bG$; moment columns are derivatives with respect to the receiver coordinate,
symmetrised with a factor $\tfrac12$ on the shear components. In that convention the contrast
operator is
\begin{equation}\label{eq:delta}
  \bD=\operatorname{blockdiag}\!\bigl(\omega^2\Delta\rho\,\bI_3,\ \Delta\mathcal{C}\bigr),
\end{equation}
where $\Delta\mathcal{C}$ is the $6\times6$ contrast stiffness with $\Delta\lambda+2\Delta\mu$ and
$\Delta\lambda$ in the normal block and $2\Delta\mu$ (not $\Delta\mu$) on the shear diagonal.

\paragraph{The laterally uniform layer.} A layer of thickness $D$ is filled with $n$ planes of
voxels, $d=D/n$, and illuminated at normal incidence. The field is then laterally uniform, only the
zero-lateral-wavenumber (specular) component of the lattice coupling acts, and
Eq.~\eqref{eq:foldylax} becomes exactly a chain of $n$ planes coupled by the specular lattice sums
\begin{equation}\label{eq:specsum}
  \bS(m)=\sum_{\mathbf R}\bG(\mathbf R+m d\,\hat{\mathbf z}),
\end{equation}
the sum running over the square lattice $\mathbf R$ of pitch $d$, $m$ the plane separation in
pitches. The question is whether, and how fast, this chain converges to the continuous layer as
$n\to\infty$ at fixed $D$.

\section{The continuous layer and its thin-layer expansion}\label{sec:continuum}

At normal incidence a plane force drives a one-dimensional problem per wave type, with modulus $M$
($M_P$ for P, $\mu$ for S), wavenumber $k$ and plane Green's function
$g(z)=\tfrac{i}{2Mk}e^{ik|z|}$. The layer's reflection and transmission follow from continuity of
displacement and traction at its faces; they agree with the stratified propagator used as the
reference in this work to $2.5\times10^{-8}$, P and S (\nb{ContinuumLimit\_Reference.wl}).

For a layer of thickness $D$ with modulus $M_1$ and density $\rho_1$ in a background $M_0,\rho_0$,
the reflection coefficient expands as
\begin{equation}\label{eq:thin}
  R=\frac{i\omega D}{2}\sqrt{\frac{M_0}{\rho_0}}
  \Bigl[\frac{\rho_1-\rho_0}{M_0}-\rho_0\Bigl(\frac1{M_1}-\frac1{M_0}\Bigr)\Bigr]+O(D^2),
\end{equation}
exact in the contrast (\nb{ContinuumLimit\_ThinLayer.wl}): to first order a thin layer is an
interface across which the traction jumps by the excess inertia and the displacement by the excess
compliance. At oblique incidence the same statement holds with the first-order system
$\partial_z\mathbf b=\bA\mathbf b$ for the P--SV state $\mathbf b=(u_x,u_z,t_{xz},t_{zz})$: the layer
acts to $O(D)$ as the jump $\bI+(\bA_1-\bA_0)D$, which reproduces the layer's $O(D)$ reflection
coefficients $R_{PP},R_{PS}$ to $10^{-12}$; the series truncated at $O(D^k)$ converges at exactly
order $k$; and a soft thin layer with shear modulus $\mu_1=D/\eta$ at fixed compliance $\eta$ tends to
the linear-slip interface as $D\to0$. These expansions are the target a single plane of voxels must
reproduce order by order.

\section{The specular sums of the point Green's tensor}\label{sec:point}

By Poisson summation,
\begin{equation}\label{eq:poisson}
  \bS(m)=\frac1{d^2}\sum_{\mathbf g}\hat\bG(\mathbf g;\,m d),\qquad \mathbf g=\tfrac{2\pi}{d}(p,q),
\end{equation}
with $\hat\bG$ the two-dimensional transform of the Green's tensor between planes. The $\mathbf g=0$
term is the continuum's plane-wave term; the others are evanescent, decaying as
$e^{-2\pi m\sqrt{p^2+q^2}}$ independently of $d$. In the strain-from-moment block they carry two
powers of $|\mathbf g|\propto1/d$, so their contribution per voxel volume does \emph{not} vanish as
$d\to0$: at $m=1$ it is eleven times the continuum term (\nb{ContinuumLimit\_SpecularSums.wl},
matched to the package's Ewald lattice sum to $5\times10^{-13}$). This surviving term is static.
Summed over all planes in slab order, the static lattice sum of the point kernel is exactly the
depolarisation of a uniform plate --- $-1/M_P$ on $\varepsilon_{zz}\leftarrow M_{zz}$ and $-1/(2\mu)$
on the normal shears, per unit volume --- plus a cubic-symmetric remainder
(\nb{ContinuumLimit\_StaticLocal.wl}, to $2\times10^{-10}$). The point kernel therefore reproduces
the continuum only if that remainder is cancelled by the voxel's own self term.

\section{The tiling identity}\label{sec:tiling}

Replace the point coupling by the \emph{source-cell average}
$\langle\bG\rangle(\mathbf r)=d^{-3}\int_{\text{cube}}\bG(\mathbf r-\mathbf y)\,d\mathbf y$ --- the
field at a voxel centre of a uniform polarisation of the source voxel. Each Poisson order in
Eq.~\eqref{eq:poisson} is then multiplied by the cube's form factor
\begin{equation}
  F(\mathbf g)=\sinc(g_x h)\,\sinc(g_y h),
\end{equation}
which vanishes at every $\mathbf g\neq0$ of the reciprocal lattice since $\sinc(p\pi)=0$ for integer
$p\neq0$. This is the Fourier form of the fact that cubes tile the plane,
$\sum_{\mathbf R}\chi(\mathbf x-\mathbf R)=1$ for the cube's indicator $\chi$. Two exact statements
follow.

\paragraph{Between planes.} For $m\neq0$ every evanescent order is removed and
\begin{equation}\label{eq:between}
  \bS_{\text{avg}}(m)=\frac1{d^2}\sum_{W=P,S}\hat\bG_W(0;\,m d)\,\sinc(k_W h),
\end{equation}
the continuum's plane-wave term times the vertical form factor. The package's cell-averaged kernel
agrees in all 81 entries to $2.8\times10^{-16}$ (\nb{ContinuumLimit\_AveragedSums.wl}).

\paragraph{Within a plane.} The sum over \emph{all} cells of the plane, the voxel's own included,
is the field at the mid-plane of a uniform plate of thickness $d$:
\begin{equation}\label{eq:tiling}
  d^3\,\bS_{\text{avg}}(0)+\bS_{\text{self}}=d^3\,\bP_0,
  \qquad
  \bS_{\text{self}}=\int_{\text{cube}}\bG(-\mathbf y)\,d\mathbf y,
\end{equation}
where $\bP_0=d^{-3}\int_{-h}^{h}\bG_{\text{plate}}(-z')\,dz'$ is elementary: in the one-dimensional
transform each entry is $c_0+A/(Mk_z^2-\rho\omega^2)$, so it averages to
$c_0/d+A\,(e^{ikh}-1)/(dMk^2)$. Its only nonzero entries are $u\leftarrow f$ on the diagonal and the
three normal strain-from-moment entries, whose local weights $-1/M_P$ and $-1/(2\mu)$ are the plate
depolarisation of \S\ref{sec:point}. Statically the identity holds for all 21 independent components
of the strain-from-moment tensor to $4\times10^{-9}$, with every cell integral evaluated analytically
(below).

\paragraph{Consequence: an exact kernel with no lattice sum.} At normal incidence the same-plane
block of the cell-averaged Foldy--Lax kernel is
\begin{equation}\label{eq:exactkernel}
  \bS_{\text{avg}}(0)=\bP_0-\bS_{\text{self}}/d^3 ,
\end{equation}
with no Ewald splitting, near shell or asymptotic tail. The route it replaces --- Ewald summation
with a near shell averaged by quadrature and an $O(d^2)$ Taylor tail beyond it --- carries a static
strain-from-moment error of $9.4\times10^{-4}$ at its default settings, falling to $2.6\times10^{-5}$
only when both the shell radius and the quadrature order are raised; Eq.~\eqref{eq:exactkernel} is
now the package's kernel at normal incidence.

\paragraph{Cell integrals without quadrature.} Write the static (Kelvin) tensor as
$G_{bc}=(c_A+c_B)\,\delta_{bc}/r-c_B\,\partial_b\partial_c r$ with
$c_B=(1/\mu-1/M_P)/8\pi$ and $c_A=1/(4\pi\mu)-c_B$. Every cell integral of its derivatives is a cell
integral of derivatives of $1/r$ and $r$. One application of the divergence theorem moves it onto
faces $x_p=c_p\pm h$, which never contain the origin, and there the two-dimensional antiderivative in
the face coordinates is continuous, so the face integral is a sum over the face's corners. The
distributional term at $r=0$ is carried by the divergence theorem itself. (A triple antiderivative
evaluated at the cube's corners is \emph{not} usable for repeated derivatives: its arctangent terms
jump across the coordinate planes through the origin, and differentiating twice there discards
exactly that term.) The self term so obtained reproduces the classical closed form of the cube's
depolarisation to machine precision.

\section{The single site}\label{sec:single}

The collocation closure takes the internal field of a voxel as uniform. The voxel's own field at its
centre is then $\bS_{\text{self}}\bD\,e$ for internal state $e$, so
$e=\psi+\bS_{\text{self}}\bD e$ and
\begin{equation}\label{eq:closure}
  \bT=d^3\,\bD\bigl(\bI-\bS_{\text{self}}\bD\bigr)^{-1}.
\end{equation}
With the kernel \eqref{eq:exactkernel} the self term cancels. Substituting
$\psi=(\bI-\bS_{\text{self}}\bD)e$ into Eq.~\eqref{eq:foldylax},
\begin{equation}\label{eq:collocation}
  e_i=\psi^0_i+d^3\sum_j \bK_{\text{all}}(i,j)\,\bD\,e_j,
  \qquad
  \bK_{\text{all}}(i,i)=\bP_0,\quad \bK_{\text{all}}(i,j)=\bS_{\text{avg}}(i-j),
\end{equation}
an identity checked on the matrices to $4\times10^{-16}$. Every block of $\bK_{\text{all}}$ is an
\emph{exact} cell integral of the one-dimensional continuum kernel (Eqs.~\eqref{eq:between},
\eqref{eq:tiling}), so Eq.~\eqref{eq:collocation} is precisely the collocation of the continuum
Lippmann--Schwinger equation with a strain piecewise constant over voxels. It does not contain
$\bS_{\text{self}}$ at all.

The package's cube $T$-matrix is this closure statically: both blocks agree with
Eq.~\eqref{eq:closure} to $2\times10^{-5}$ at $k_Sa=0.01$ ($a$ the half-side), the difference falling
as $(ka)^2$ (\nb{ContinuumLimit\_CubeT.wl}). It departs dynamically by about $0.25(ka)^2$ through its
far-field form-factor corrections. Its self integrals are exact in the static and radiation parts
but omit the real even-power terms of the Green's tensor's expansion ($k^2r$, $r^3$, \dots), an
error of $0.4$--$0.6\,(ka)^2$ against the analytic self term, whose series agrees with the
closed-form elastodynamic tensor to $10^{-14}$. By Eq.~\eqref{eq:collocation} that error cannot enter
the closure scheme.

\section{Convergence, and the closed-form error}\label{sec:convergence}

\paragraph{Second order.} For a layer $D=2$\,m with the contrast
$\Delta\lambda=2$\,GPa, $\Delta\mu=1$\,GPa, $\Delta\rho=100$\,kg\,m$^{-3}$ in a background
$\alpha=5000$, $\beta=3000$\,m\,s$^{-1}$, $\rho=2500$\,kg\,m$^{-3}$, illuminated by a P plane force,
the reflected displacement of the chain converges to the continuous layer's at order $2.000$ for
$n=4\to16$ at $\omega=60$, $300$ and $600$\,rad\,s$^{-1}$, with the collocation
\eqref{eq:collocation} and with the package $T$-matrix alike (\nb{ContinuumLimit\_Chain.wl}). The
error depends on $k_Sa$ alone.

\paragraph{The leading error in closed form.} In Eq.~\eqref{eq:collocation} the kernel is
integrated exactly; the only approximation is the internal field's variation across a voxel. At
first order in the contrast the discrete scattered field is therefore a midpoint rule applied to
the continuum's Born integral, with the kernel's cell integral exact. For reflection the integrand
carries the two-way phase $e^{2ikz}$; for transmission it is flat. Hence, for \emph{every} $n$,
\begin{equation}\label{eq:errorR}
  \frac{R_{\text{disc}}}{R}=\frac{\sinc(kh)}{\sinc(2kh)}=1+\frac{(kd)^2}{8}+O(d^4),
  \qquad
  \frac{T^{\text{sc}}_{\text{disc}}}{T^{\text{sc}}}=\sinc(kh)=1-\frac{(kd)^2}{24}+O(d^4),
\end{equation}
confirmed for $n=1$ to $16$ to $10^{-4}$ of the error itself
(\nb{ContinuumLimit\_ErrorConstant.wl}). At the full contrast the reflection constant is unchanged
($0.124$ against $1/8$); the transmission constant grows to $-0.092$.

\paragraph{Why a uniform-field single site stops at second order.} The two constants have opposite
signs, so no scalar correction of $\bT$ --- no single $O((ka)^2)$ factor --- can cancel both. The
package $T$-matrix illustrates it: its form-factor corrections reduce the reflection error by a
factor $4.45$ and \emph{increase} the transmission error by $1.7$, at every $n$. The missing term in
Eq.~\eqref{eq:errorR} is the first moment of the internal field over each voxel, which radiates
through the derivative of the kernel. A fourth-order scheme must carry it: a single site with a
gradient channel, and a kernel whose cell integrals include the kernel's derivative.

\section{Fourth order: the first-moment voxel}\label{sec:fourth}

Let each voxel carry the mean \emph{and} the first moment of its internal field. On voxel $j$, with
local coordinate $s\in[-h,h]$ and Legendre functions $\phi_0=1$, $\phi_1=s/h$, the field is
$w=\sum_{b}c_{b,j}\,\phi_b$, $w=(u,\varepsilon)$ the displacement and strain of the P channel, and
the polarisation is $q=(\omega^2\Delta\rho\,u,\ \Delta M\,\varepsilon)$ with $\Delta M=\Delta\lambda
+2\Delta\mu$. Testing the continuum equation
\begin{equation}
  w(z)=w^0(z)+\int\bK(z-z')\,q(z')\,dz',\qquad
  \bK=\begin{pmatrix} g & g' \\ g' & g''\end{pmatrix},\quad g''=-k^2g-\delta(z)/M_P,
\end{equation}
with the same two functions on every voxel gives
\begin{equation}\label{eq:galerkin}
  \Bigl(\int_{i}\phi_a^2\Bigr)c_{a,i}
  =\int_i\phi_a w^0+\sum_{j,b}\Bigl[\int_i\!\!\int_j\phi_a(z)\,\bK(z-z')\,\phi_b(z')\,dz\,dz'\Bigr]
  \bD_1 c_{b,j},
\end{equation}
$\bD_1=\operatorname{diag}(\omega^2\Delta\rho,\Delta M)$. Between voxels the double integral is
separable, a product of Legendre moments of $e^{\pm ikz}$; within a voxel it is split at $z'=z$ and
the local term contributes $-\int_i\phi_a\phi_b/M_P$. Every entry is closed-form. Keeping only
$\phi_0$ recovers a second-order scheme; keeping $\phi_0$ and $\phi_1$ converges at fourth order in
both reflection and transmission (\nb{ContinuumLimit\_FourthOrder.wl}):
\begin{center}
\begin{tabular}{@{}rccc@{}}
\toprule
$n$ & collocation & mean only & mean and first moment \\
\midrule
$1$  & $1.8\times10^{-3}$ / $1.3\times10^{-3}$ & $1.2\times10^{-3}$ / $1.8\times10^{-3}$ &
       $2.8\times10^{-7}$ / $2.8\times10^{-7}$ \\
$4$  & $1.1\times10^{-4}$ / $8.3\times10^{-5}$ & $7.5\times10^{-5}$ / $1.1\times10^{-4}$ &
       $1.1\times10^{-9}$ / $1.1\times10^{-9}$ \\
$16$ & $7.0\times10^{-6}$ / $5.2\times10^{-6}$ & $4.7\times10^{-6}$ / $6.9\times10^{-6}$ &
       $4.2\times10^{-12}$ / $4.2\times10^{-12}$ \\
order & $2$ / $2$ & $2$ / $2$ & $4$ / $4$ \\
\bottomrule
\end{tabular}
\end{center}
(relative error of the scattered displacement, reflection / transmission, $\omega=300$\,rad\,s$^{-1}$;
the order is $4.00$ at $\omega=60$ and $600$\,rad\,s$^{-1}$ too). A single plane of first-moment
voxels represents the $2$\,m layer $6500$ times more accurately than a single plane of uniform ones.
The collocation column is the scheme of \S\ref{sec:single} and reproduces the chain of the full
$9\times9$ system to all printed digits.

\paragraph{In three dimensions.} The three-dimensional first-moment voxel carries the nine-component
state with Legendre weights $\{1,\,x/h,\,y/h,\,z/h\}$ --- 36 unknowns --- tested with the same
functions and coupled by cell-to-cell double integrals of the Green's tensor. For a laterally uniform
layer at normal incidence it reduces \emph{exactly} to Eq.~\eqref{eq:galerkin}
(\nb{ContinuumLimit\_Reduction3D.wl}). The lateral first moments are never excited: the lattice and
the incident field are symmetric under reflection about every voxel's centre planes, and $x/h$, $y/h$
are odd there. The mean and the $z$-moment carry weights that are constant laterally, so their lateral
form factor $\sinc(g_xh)\sinc(g_yh)$ vanishes on the reciprocal lattice and the plane sum of the
double integrals is the plate-to-plate integral in $z$ alone. Two facts mark where genuinely
three-dimensional structure begins. The lateral first moment's form factor,
$F_1(g)=i\,(gh\cos gh-\sin gh)/(gh)^2$, does not vanish on the reciprocal lattice
($F_1=i(-1)^p/p\pi$), so the gradient channels couple through the evanescent orders. And at oblique
incidence the mean's own zeros shift, $\sinc\bigl((2\pi p/d+k_x)h\bigr)=(-1)^p k_xh/p\pi+O(k_x^2h^2)$,
so the tiling identity holds only to $O(k_\parallel h)$.

\paragraph{Oblique incidence.} The scheme is built spectrally. By Poisson summation the Bloch coupling
between planes $i$ and $j$ at lateral wavenumber $\mathbf k_\parallel$ is
\begin{equation}\label{eq:bloch}
  C_{ab}(i,j)=\frac1{d^2}\sum_{\mathbf g}L_a(\boldsymbol\kappa)\,L_b(-\boldsymbol\kappa)\,
  \int\!\!\int\phi_a(z)\,\hat\bG(\boldsymbol\kappa;z-z')\,\phi_b(z')\,dz\,dz',
  \qquad \boldsymbol\kappa=\mathbf k_\parallel+\mathbf g,
\end{equation}
with $L_a$ the lateral Legendre factors. The whole-space tensor decays as $1/k_z^2$ and the nine-component
state adds at most two powers of $k_z$, so each entry of $\hat\bG(\boldsymbol\kappa;z)$ is a $\delta(z)$
of constant weight plus, per mode, $[i\alpha/2\gamma+i\beta\,\mathrm{sgn}(z)/2]\,e^{i\gamma|z|}$,
$\gamma^2=k_W^2-\kappa^2$: every $z$ integral is closed-form. Double averaging makes the lattice sum
converge without Ewald splitting, to three digits at $|p|,|q|\le8$. Against the exact layer, for a P
wave incident at $20^\circ$ and $40^\circ$ (\nb{ContinuumLimit\_Oblique.wl}):
\begin{center}
\small\setlength{\tabcolsep}{4pt}
\begin{tabular}{@{}rcccc@{}}
\toprule
 & \multicolumn{2}{c}{mean only} & \multicolumn{2}{c}{mean and first moments}\\
$n$ & $R_{PP}$ & $R_{PS}$ & $R_{PP}$ & $R_{PS}$ \\
\midrule
$1$ & $9.3\times10^{-4}$ & $1.7\times10^{-3}$ & $7.2\times10^{-8}$ & $4.1\times10^{-7}$ \\
$2$ & $2.3\times10^{-4}$ & $4.2\times10^{-4}$ & $4.5\times10^{-9}$ & $2.5\times10^{-8}$ \\
$8$ & $1.5\times10^{-5}$ & $2.7\times10^{-5}$ & $1.7\times10^{-11}$ & $9.9\times10^{-11}$ \\
order & $2$ & $2$ & $4$ & $4$ \\
\bottomrule
\end{tabular}
\end{center}
(relative error, $20^\circ$; at $40^\circ$ the orders are again $2$ and $4$). The mean-only voxel is
second order and the first-moment voxel fourth order in both the reflected P and the converted S wave:
the breakdown of the tiling identity at $O(k_\parallel h)$ and the evanescent coupling of the gradient
channels do not cost an order. At $\mathbf k_\parallel=0$ the construction reproduces
Eq.~\eqref{eq:galerkin} to all printed digits.

\section{Summary}\label{sec:summary}

\begin{center}
\begin{tabular}{@{}>{\raggedright\arraybackslash}p{0.52\linewidth}>{\raggedright\arraybackslash}p{0.40\linewidth}@{}}
\toprule
Result & Evidence \\
\midrule
Continuous layer, exact and to $O(D^k)$ & reference to $2.5\times10^{-8}$; series order $k$ \\
Point-kernel local term survives $d\to0$ & $11\times$ the continuum term at $m=1$ \\
Averaged coupling between planes $=$ continuum$\times\sinc$ & $2.8\times10^{-16}$, 81 entries \\
Tiling identity, Eq.~\eqref{eq:tiling} & 21/21 components to $4\times10^{-9}$ \\
Exact kernel, Eq.~\eqref{eq:exactkernel} & replaces a $9.4\times10^{-4}$ static error \\
Self term cancels, Eq.~\eqref{eq:collocation} & $4\times10^{-16}$ \\
Package $T$ $=$ closure statically & $2\times10^{-5}$ at $k_Sa=0.01$ \\
Second-order convergence & order $2.000$, three frequencies \\
Closed-form error, Eq.~\eqref{eq:errorR} & $n=1$--$16$, to $10^{-4}$ of the error \\
First-moment voxel, Eq.~\eqref{eq:galerkin} & order $4.00$, three frequencies \\
First-moment voxel, oblique, Eq.~\eqref{eq:bloch} & order $4$, P and S, $20^\circ$ and $40^\circ$ \\
\bottomrule
\end{tabular}
\end{center}

The discrete scheme is exact wherever the voxel is: its coupling is the continuum's cell integral,
and its self term cancels. What remains is the assumption that the field inside a voxel is uniform,
and that error is known in closed form. A voxel that also carries the first moment of its field
removes it and converges at fourth order: in one dimension, exactly in three at normal incidence,
and in three at oblique incidence for both the reflected P and the converted S wave.

\vfill
\noindent\rule{\linewidth}{0.4pt}\\
{\footnotesize Every numerical value is reproducible from the Mathematica notebooks
\path{Mathematica/ContinuumLimit_X.wl}, for X in Reference, ThinLayer, SpecularSums,
StaticLocal, AveragedSums, CubeT, Chain, ErrorConstant, FourthOrder, Reduction3D and Oblique, and their reference dumps
\path{scripts/continuum_limit_*.py}; the exact kernel is
\path{cubic_scattering/cell_averaged_lattice.py}.}

\end{document}
